% Appendix

\section{Page-Vector Encoding: Formal Specification}
\label{app:pvs-encoding}

This appendix formalizes the page-vector construction used by the UDEG
representation.  The purpose of the encoding is to turn a mobile observation
into a compact state key that supports fast lookup and action relocation,
without storing raw screenshots, raw XML, or source-screen text in the shared
payload.

\paragraph{Element vector.}
Let $e_i$ denote the $i$-th element in an accessibility tree with $n$ elements.
OmniFlow maps $e_i$ to a 64-dimensional vector
$\phi(e_i) \in \mathbb{R}^{64}$, formed by concatenating two feature blocks:
\[
\begin{aligned}
  h(e_i) &= [c_i;t_i;a_i;m_i;s_i] \in
  \mathbb{R}^{16+4+4+4+4},\\
  g(e_i) &= [k_i;b_i;u_i;\ell_i;d_i] \in
  \mathbb{R}^{8+8+12+2+2}.
\end{aligned}
\]
Here $c_i$ is a text or visual-content hash, $t_i$ is content type, $a_i$ is
affordance, $m_i$ is visual prominence, $s_i$ is visual state, $k_i$ is widget
class, $b_i$ is boolean attributes, $u_i$ is bounded subtree structure,
$\ell_i$ is hierarchy position, and $d_i$ is the resource-id hint.  The final
embedding is:
\[
  \phi(e_i)=
  [\lambda_h\,\nu(h(e_i)) \,;\, \lambda_p\,\nu(g(e_i))],
  \quad
  \nu(\mathbf{x})=
  \begin{cases}
    \mathbf{x}/\|\mathbf{x}\|_2, & \|\mathbf{x}\|_2>0,\\
    \mathbf{0}, & \text{otherwise}.
  \end{cases}
\]
The current implementation uses $\lambda_h=0.55$ and $\lambda_p=0.45$.

\emph{Human-view block.}
The human-view block $h(e_i) \in \mathbb{R}^{32}$ encodes rendered-UI cues:
a 16-dimensional content hash, 4 content-type dimensions, 4 affordance
dimensions, 4 visual-prominence dimensions, and 4 visual-state dimensions.

\emph{Programmer-view block.}
The programmer-view block $g(e_i) \in \mathbb{R}^{32}$ encodes
accessibility-XML cues: an 8-dimensional widget-class hash, 8 interaction and
state flags, a 12-dimensional bounded-subtree hash, 2 hierarchy dimensions,
and 2 resource-id hint dimensions.  Raw \texttt{text} and
\texttt{content-desc} strings are not serialized in the minimal shared
payload.  Local XML-built pages may keep hashed text and selector diagnostics
for action mapping, but these fields are opt-in diagnostics rather than the
default shared page-vector payload.

\paragraph{Slice pooling.}
OmniFlow pools element vectors into eight page slices.  The current slice order
is:
$R_1$ = global text,
$R_2$ = top-bar text,
$R_3$ = footer/navigation text,
$R_4$ = action text,
$R_5$ = state text,
$R_6$ = global UI structure,
$R_7$ = body structure, and
$R_8$ = front-surface layout.
The top-bar and footer cuts are quantile-based rather than fixed 10\% screen
bands, so the same encoder can adapt to pages with different root heights and
system UI insets.

For each slice $R_r$, the page-level vector is pooled with state-aware element
weights:
\[
  \bar{\mathbf{p}}_r =
  \frac{\sum_{e_i \in E_r}\omega_i\phi(e_i)}
       {\sum_{e_i \in E_r}\omega_i},
\]
where $E_r$ is the set of elements assigned to the slice and $\omega_i$ is a
positive element weight.  If $E_r$ is empty, $\bar{\mathbf{p}}_r =
\mathbf{0}$.  Each nonzero slice is normalized to $\hat{\mathbf{p}}_r$.

The page-level recall vector is the slice-weighted concatenation:
\[
  \mathbf{p} =
  [w_1\hat{\mathbf{p}}_1;\ldots;w_8\hat{\mathbf{p}}_8]
  \in \mathbb{R}^{512},
\]
with fixed slice weights $w_r$.  The vector store normalizes this final vector
to unit length for cosine lookup.

\begin{algorithm}[t]
\caption{XML-to-page vector encoding}
\label{alg:xml-page-embedding}
\small
\begin{algorithmic}[1]
\Require XML observation $X$ and optional screenshot $I$
\Ensure normalized page embedding $\mathbf{p}$
\State $E \gets \textsc{ParseElements}(X)$
\ForAll{$e_i \in E$}
  \State $\mathbf{h}_i \gets \textsc{HumanView}(e_i, I)$
  \State $\mathbf{g}_i \gets \textsc{ProgrammerView}(e_i, X)$
  \State $\phi(e_i) \gets [\lambda_h\nu(\mathbf{h}_i);\lambda_p\nu(\mathbf{g}_i)]$
  \State $\omega_i \gets \textsc{ElementWeight}(e_i, X)$
\EndFor
\For{$r = 1$ to $8$}
  \State $E_r \gets \{e_i \in E \mid \textsc{Slice}(e_i, X)=r\}$
  \If{$E_r = \emptyset$}
    \State $\mathbf{p}_r \gets \mathbf{0}$
  \Else
    \State $\bar{\mathbf{p}}_r \gets \frac{\sum_{e_i\in E_r}\omega_i\phi(e_i)}{\sum_{e_i\in E_r}\omega_i}$
    \State $\mathbf{p}_r \gets \textsc{NormalizeSlice}(\bar{\mathbf{p}}_r)$
  \EndIf
\EndFor
\State $\mathbf{p} \gets [w_1\mathbf{p}_1;\ldots;w_8\mathbf{p}_8]$
\State $\mathbf{p} \gets \textsc{NormalizeForLookup}(\mathbf{p})$
\State \Return $\mathbf{p}$
\end{algorithmic}
\end{algorithm}

\paragraph{Page similarity.}
The similarity between two pages is computed over normalized page-level
vectors:
\[
  \mathrm{sim}(\mathbf{p}^{(1)},\mathbf{p}^{(2)})
  = \mathbf{p}^{(1)} \cdot \mathbf{p}^{(2)}.
\]
This fixed 512-dimensional vector is used for fast candidate recall.  The
serialized page-vector payload used by action transfer is larger because it
also contains element vector records, root bounds, and element bounds.  The
evaluation therefore reports the actual recalled payload size emitted by the
runner instead of assuming a constant page-vector size.

\paragraph{Privacy invariant.}
The encoding pipeline satisfies the following operational invariant:
\emph{the minimal shared page-vector payload contains no raw text or raw
content-description string from the source accessibility tree}.  Element
records retain bounds because action relocation needs geometry.  Local
diagnostic payloads may additionally include resource IDs, class names, flags,
and text/content-description hashes when action mapping needs exact signature
evidence.  The default shared serialization omits those identity-bearing
fields.  This is an operational privacy boundary, not a formal
differential-privacy claim.

\section{On-Device Observation Description Cost}
\label{app:on-device-cost}

Table~\ref{tab:appendix-on-device-cost} records the appendix-only diagnostic
snapshot used to contextualize on-device memory lookup cost.  The snapshot
measures the time and token scale for converting one mobile observation into a
complete textual description.  It is not an end-to-end task result, a success
rate, or an energy measurement.  Its role is to show why repeated
screenshot/XML description is a bottleneck for an on-device replay system, and
why OmniFlow uses PageVector lookup as a 0-token state key.

\begin{center}
  \refstepcounter{table}\label{tab:appendix-on-device-cost}
  \small\textbf{Table~\thetable.} Appendix-only single-observation description
  cost snapshot.
  \vspace{0.4em}

  \resizebox{\columnwidth}{!}{%
  \begin{tabular}{llllr}
    \toprule
    Location & Backend & Input & Tokens & Time \\
    \midrule
    On-device & Llama.cpp + 0.8B-W4 + CPU & screenshot & $\sim$2,500 & 60s+ \\
    On-device & MNN-LMM + 0.8B-W4 + CPU & screenshot & $\sim$2,500 & 10s+ \\
    On-device & Omniinfer + 0.8B-W4 + NPU & XML & $\sim$1k--10k+ & 3s \\
    Cloud & Qwen-VL-Plus & screenshot & $\sim$3,200 & 3s \\
    \midrule
    OmniFlow & PageVector & XML features & 0 & 4.13 ms/item \\
    \bottomrule
  \end{tabular}}
\end{center}

\section{AITW Prompt Similarity Characterization}
\label{app:aitw-prompt-similarity}

\paragraph{Dataset and denominator.}
We use the prompt vocabulary of Android in the Wild (AITW)~\cite{rawles2023android}.
AITW reports unique-prompt counts separately for General, GoogleApps, Install,
Single, and WebShopping.  We preserve this official five-subset accounting, so
the denominator is the sum of the five counts, 30,378.  Merging identical raw
strings across subsets yields 30,165 globally distinct strings; the difference
comes from 213 additional cross-subset occurrences.  We extracted the strings
through a processed public index\footnote{\url{https://huggingface.co/datasets/KMK040412/aitw-processed-labeled-full}},
which covers 5,687,349 interaction steps and 714,327 episodes.  Its prompt
vocabulary matches the official count in every subset.  We form candidate
peers over the global prompt pool and project the resulting indicator back to
the official subset-level entries.

\paragraph{Parameter templates.}
We canonicalize case and punctuation, then replace URLs, email addresses,
quoted values, and numeric values with typed placeholders.  This operation
groups instructions that express the same operation with different task
arguments.  A prompt is counted by the template-only statistic when another
prompt maps to the same canonical template.

\paragraph{Semantic matching.}
We encode each globally distinct prompt with
\texttt{BAAI/bge-small-en-v1.5} and compute cosine similarity between normalized
embeddings.  For threshold $\tau$, prompt $p$ is marked as having a reuse
candidate when at least one other prompt $q$ satisfies
\[
  M_\tau(p)=\mathbf{1}\!\left[
  \begin{aligned}
    \exists q\neq p:\quad &T(p)=T(q),\quad\text{or}\\
    &\cos(z_p,z_q)\geq\tau \land \neg C(p,q)
  \end{aligned}
  \right],
\]
where $T(\cdot)$ is the parameter template, $z$ is the prompt embedding, and
$C$ is an explicit contradiction filter.  The filter rejects pairs that
reverse the requested operation, including on versus off, enable versus
disable, install versus uninstall, add versus remove or delete, create versus
delete, open versus close, show versus hide, login versus logout, subscribe
versus unsubscribe, allow versus block, accept versus reject, and increase
versus decrease.  We report
\[
  R_\tau=\frac{1}{30{,}378}\sum_p M_\tau(p).
\]
The main analysis uses $\tau=0.98$; $\tau=0.95$ is reported only as a
sensitivity point.

\begin{table}[t]
  \centering
  \caption{AITW prompt-level reuse-candidate estimate.  The main paper uses the
  strict $\tau=0.98$ result.}
  \label{tab:aitw-prompt-sensitivity}
  \small
  \begin{tabular}{lrr}
    \toprule
    Criterion & Similar / Total & Rate \\
    \midrule
    Parameter template only & 11,052 / 30,378 & 36.4\% \\
    Template or cosine $\geq 0.98$ & 22,834 / 30,378 & 75.2\% \\
    Template or cosine $\geq 0.95$ & 26,386 / 30,378 & 86.9\% \\
    \bottomrule
  \end{tabular}
\end{table}

\begin{table}[t]
  \centering
  \caption{Strict AITW result by official subset ($\tau=0.98$).}
  \label{tab:aitw-prompt-subsets}
  \small
  \begin{tabular}{lrr}
    \toprule
    Subset & Similar / Total & Rate \\
    \midrule
    General & 69 / 545 & 12.7\% \\
    GoogleApps & 81 / 306 & 26.5\% \\
    Install & 549 / 688 & 79.8\% \\
    Single & 9,151 / 15,366 & 59.6\% \\
    WebShopping & 12,984 / 13,473 & 96.4\% \\
    \bottomrule
  \end{tabular}
\end{table}

\paragraph{Scope.}
The unit above is a prompt, so $R_\tau$ is a Function-reuse candidate upper
bound rather than an executable-reuse measurement.  The analysis does not
establish that paired episodes share a compatible start state, action sequence,
argument binding, or terminal effect.

\paragraph{Trace test.}
A separate AITW trace-correspondence protocol can sample distinct source and
target episodes from the same app and
parameter template, add same-app near-miss and opposite-intent hard negatives,
and split by held-out templates.  Function Recall@$k$, segment-alignment
accuracy, and abstain rate then measure offline generalization.  End-to-end
replay success still requires a live environment with control-aware execution
and terminal/effect verification.

% % Appendix
%
% \appendix
%
% \section{Page-Vector Encoding: Formal Specification}
% \label{app:pvs-encoding}
%
% This appendix formalises the page-vector construction described informally
% in \S\ref{sec:design-offline}.  The goal is to show precisely what
% information is captured, what is discarded, and why the representation remains
% available when semantic attributes are absent while preserving the paper's
% privacy boundary (no raw text or raw screenshots are retained).
%
% \paragraph{Element vector.}
% Let $e_i$ denote the $i$-th element in an accessibility tree with $n$
% elements.  OmniFlow maps $e_i$ to a 64-dimensional real vector
% $\phi(e_i) \in \mathbb{R}^{64}$, formed by concatenating the two
% feature blocks in \texttt{ElementEmbeddingConfig}.  The human-view block
% encodes visible cues, while the programmer-view block encodes accessibility
% XML structure and execution cues:
% \[
% \begin{aligned}
%   h(e_i) &= [c_i;t_i;a_i;m_i;s_i] \in
%   \mathbb{R}^{16+4+4+4+4},\\
%   g(e_i) &= [k_i;b_i;u_i;\ell_i;d_i] \in
%   \mathbb{R}^{8+8+12+2+2},
% \end{aligned}
% \]
% where $c_i$ is the text or visual-content hash, $t_i$ is content type,
% $a_i$ is affordance, $m_i$ is visual prominence, $s_i$ is visual state,
% $k_i$ is widget class, $b_i$ is boolean attributes, $u_i$ is bounded subtree
% structure, $\ell_i$ is hierarchy position, and $d_i$ is the resource-id hint.
% The final element embedding is:
% \[
%   \phi(e_i)=
%   [\lambda_h\,\nu(h(e_i)) \,;\, \lambda_p\,\nu(g(e_i))],
%   \quad
%   \nu(\mathbf{x})=
%   \begin{cases}
%     \mathbf{x}/\|\mathbf{x}\|_2, & \|\mathbf{x}\|_2>0,\\
%     \mathbf{0}, & \text{otherwise}.
%   \end{cases}
% \]
% The current implementation uses $\lambda_h=0.55$ and $\lambda_p=0.45$.
%
% \emph{Human-view block} $h(e_i) \in \mathbb{R}^{32}$.
% This block encodes rendered-UI cues: a 16-dimensional content hash, 4
% content-type dimensions, 4 affordance dimensions, 4 visual-prominence
% dimensions, and 4 visual-state dimensions.
%
% \emph{Programmer-view block} $g(e_i) \in \mathbb{R}^{32}$.
% This block encodes accessibility-XML cues: an 8-dimensional widget-class hash,
% 8 interaction/state flags, a 12-dimensional bounded-subtree hash, 2 hierarchy
% dimensions, and 2 resource-id hint dimensions.
% Raw \texttt{text} and \texttt{content-desc} strings are not serialized in the
% minimal shared payload.  Local XML-built pages can keep hashed text and
% selector/state diagnostics for action mapping, but those fields are opt-in
% diagnostics rather than the default shared page-vector payload.
%
% \paragraph{Slice pooling.}
% OmniFlow's page-level recall vector pools element vectors into eight slices.
% The current slice order is:
% $R_1$ = global text,
% $R_2$ = top-bar text,
% $R_3$ = footer/navigation text,
% $R_4$ = action text,
% $R_5$ = state text,
% $R_6$ = global UI structure,
% $R_7$ = body structure, and
% $R_8$ = front-surface layout.
% The top-bar and footer cuts are quantile-based rather than fixed 10\% screen
% bands, so the same encoder can adapt to pages with different root heights and
% system UI insets.
%
% The page-level vector for slice $R_r$ is pooled with state-aware element
% weights:
% \[
%   \bar{\mathbf{p}}_r =
%   \frac{\sum_{e_i \in E_r}\omega_i\phi(e_i)}
%        {\sum_{e_i \in E_r}\omega_i},
% \]
% where $E_r$ is the set of elements assigned to slice $R_r$ and $\omega_i$ is a
% positive element weight.  If $E_r$ is empty, $\bar{\mathbf{p}}_r =
% \mathbf{0}$.  Each nonzero slice is then normalized to
% $\hat{\mathbf{p}}_r$.
%
% The page-level recall vector is the slice-weighted concatenation
% \[
%   \mathbf{p} =
%   [w_1\hat{\mathbf{p}}_1;\ldots;w_8\hat{\mathbf{p}}_8]
%   \in \mathbb{R}^{512},
% \]
% with fixed slice weights $w_r$.  The vector store normalizes this final vector
% to unit length for cosine lookup.
%
% \begin{algorithm}[t]
% \caption{XML-to-page embedding}
% \label{alg:xml-page-embedding}
% \small
% \begin{algorithmic}[1]
% \Require XML observation $X$ and optional screenshot $I$
% \Ensure normalized page embedding $\mathbf{p}$
% \State $E \gets \textsc{ParseElements}(X)$
% \ForAll{$e_i \in E$}
%   \State $\mathbf{h}_i \gets \textsc{HumanView}(e_i, I)$
%   \State $\mathbf{g}_i \gets \textsc{ProgrammerView}(e_i, X)$
%   \State $\phi(e_i) \gets [\lambda_h\nu(\mathbf{h}_i);\lambda_p\nu(\mathbf{g}_i)]$
%   \State $\omega_i \gets \textsc{ElementWeight}(e_i, X)$
% \EndFor
% \For{$r = 1$ to $8$}
%   \State $E_r \gets \{e_i \in E \mid \textsc{Slice}(e_i, X)=r\}$
%   \If{$E_r = \emptyset$}
%     \State $\mathbf{p}_r \gets \mathbf{0}$
%   \Else
%     \State $\bar{\mathbf{p}}_r \gets \frac{\sum_{e_i\in E_r}\omega_i\phi(e_i)}{\sum_{e_i\in E_r}\omega_i}$
%     \State $\mathbf{p}_r \gets \textsc{NormalizeSlice}(\bar{\mathbf{p}}_r)$
%   \EndIf
% \EndFor
% \State $\mathbf{p} \gets [w_1\mathbf{p}_1;\ldots;w_8\mathbf{p}_8]$
% \State $\mathbf{p} \gets \textsc{NormalizeForLookup}(\mathbf{p})$
% \State \Return $\mathbf{p}$
% \end{algorithmic}
% \end{algorithm}
%
% \paragraph{Page similarity.}
% The similarity between two pages is computed over the normalized page-level
% vectors, so cosine similarity is the primary page-recall score:
% \[
%   \mathrm{sim}(\mathbf{p}^{(1)},\mathbf{p}^{(2)})
%   = \mathbf{p}^{(1)} \cdot \mathbf{p}^{(2)},
% \]
% after both vectors have been normalized.  This fixed 512-dimensional vector
% is used for fast candidate recall.  The serialized page-vector payload used by
% action transfer is larger: it also contains element vector records, root
% bounds, and element bounds, so its byte size still depends on the number of
% elements.  This is why the evaluation reports the actual recalled payload
% size emitted by the runner instead of assuming a constant page-vector size.
%
% \paragraph{Privacy invariant.}
% The encoding pipeline satisfies the following invariant by construction:
% \emph{the minimal shared page-vector payload contains no raw text or raw
% content-description string from the source accessibility tree}.  Element
% records retain bounds because coordinate transfer needs geometry; local
% diagnostic payloads may additionally include resource IDs, class names, flags,
% and text/content-description hashes when action mapping needs exact signature
% evidence.  The default shared serialization omits those action-identity fields.
% This is an operational privacy boundary, not a formal differential-privacy
% claim.  It also explains a current limitation recorded in the reviewer-risk
% ledger: a privacy-minimal page-vector payload is sufficient as a page-recall object,
% but exact action-signature transfer may need local XML context or an
% identity-bearing diagnostic payload and must otherwise abstain.
%
% % ──────────────────────────────────────────────────────────────────────────
% \section{MobileViews Attribute Completeness Analysis}
% \label{app:mobileviews-analysis}
%
% \S\ref{sec:obs2} states that 32.5\% of real Android touch targets carry no
% usable semantic attribute (\texttt{resource-id}, \texttt{text}, or
% \texttt{content-desc}).  This appendix describes the measurement methodology
% to make the claim reproducible and to clarify what ``usable'' means in
% practice.
%
% \paragraph{Dataset.}
% The analysis uses MobileViews~\cite{mobileviews}, a large-scale dataset of
% Android app UI states and recorded interaction traces collected from 30,000+
% apps.  We extract the 7,378 elements that are the direct targets of
% recorded touch actions---the set of elements that a replay or
% attribute-based agent must be able to identify.
%
% \paragraph{Attribute filtering.}
% For each touch-target element, we examine three attributes:
%
% \begin{enumerate}
% \item \texttt{resource-id}: Non-empty and \emph{not obfuscated}.  An
%   \texttt{resource-id} value is classified as obfuscated if its suffix
%   (the portion after the last ``\texttt{/}'') consists solely of lower-case
%   letters and digits and has length $\le 3$.  These are characteristic
%   outputs of ProGuard and R8 resource renaming; an agent matching on
%   \texttt{resource-id} = \texttt{a0} on one app version may find
%   \texttt{resource-id} = \texttt{b1} on the next build.
%
% \item \texttt{text}: Non-empty after trimming whitespace.  This includes
%   visible button labels, input-field content, and list-item text.  Dynamic
%   content (e.g., a contact name in a list item) is not excluded here because
%   even dynamic text can serve as a selector when the content is known at
%   action-generation time.
%
% \item \texttt{content-desc}: Non-empty after trimming whitespace.  This
%   attribute is the primary accessibility label for icon-only controls and
%   image buttons; its absence is common in apps that do not follow accessibility
%   best practices.
% \end{enumerate}
%
% An element is classified as \emph{attribute-absent} if all three values are
% empty or obfuscated after the filters above.  We do not count fallback
% attributes such as \texttt{hint}, \texttt{view-id-resource-name}, or inferred
% positional labels, because attribute-only replay baselines do not rely on
% them for robust cross-version matching.
%
% \paragraph{Result.}
% Of the 7,378 touch-target elements, 2,398 are attribute-absent under the
% definition above: 32.5\%.  The remaining 67.5\% carry at least one usable
% attribute.  Broken down by attribute type:
% \begin{itemize}
%   \item 3,923 (53.2\%) have a non-obfuscated \texttt{resource-id},
%   \item 2,224 (30.1\%) have non-empty \texttt{text},
%   \item 635 (8.6\%) have non-empty \texttt{content-desc}.
% \end{itemize}
% Attributes are not mutually exclusive; many elements carry two or all three.
% The 32.5\% attribute-absent set represents the \emph{structural failure region}
% of any attribute-based or selector-based approach: no algorithm that requires
% a stable attribute can address these elements, regardless of its sophistication.
%
% \paragraph{Generalisability.}
% MobileViews was collected from a broad sample of consumer Android apps, not
% from a curated or enterprise-app subset.  However, it skews toward apps that
% expose an accessible UI (the accessibility service must be able to observe
% the UI tree).  Apps built entirely on Unity, Unreal, or custom OpenGL
% rendering expose no accessibility tree at all; for these, the ``0 usable
% attributes'' rate is 100\%.  The 32.5\% figure is therefore a \emph{lower
% bound} on the structural failure rate in the wild.
%
% % ──────────────────────────────────────────────────────────────────────────
% \section{Control Plane Signal Taxonomy}
% \label{app:control-effects}
%
% This appendix lists the full set of correction signals the control plane may
% return, as referenced in \S\ref{sec:design-online}.  Each signal is a
% high-level semantic directive interpreted by the executor; no signal issues
% device commands directly.
%
% \begin{description}
%   \item[\texttt{replace\_action}.]  Substitute a corrected action for the
%     scheduled one before dispatch.  Used when a detector determines that the
%     planned action targets the wrong element or uses incorrect parameters for
%     the current device state (e.g., coordinate drift detected by the anchor
%     failure detector).
%
%   \item[\texttt{run\_actions}.]  Inject a short recovery action sequence
%     before resuming the Function plan at the current step index.  Used to
%     handle transient interruptions such as permission dialogs, notification
%     banners, or loading overlays that must be dismissed before the planned
%     action can proceed.
%
%   \item[\texttt{refresh\_observation}.]  Re-observe the current page without
%     advancing the plan.  Used when the page state is expected to change
%     shortly (e.g., a loading spinner is present) and the next planned action
%     should not be dispatched until the page has stabilised.
%
%   \item[\texttt{abort}.]  Stop the Function and return a structured failure
%     to the caller.  Used when a detector determines that the execution cannot
%     be recovered locally and the deviation does not warrant a host-agent
%     handoff (e.g., the package identity is irrecoverably wrong).
%
%   \item[\texttt{handoff}.]  Suspend the Function and pass the accumulated
%     execution context---completed steps as typed action records, the current
%     page vector, and the remaining plan---to the host agent.
%     This is the escalation path for deviations that require open-ended
%     reasoning.  The agent receives a partially-completed trajectory, not a
%     cold-start observation, and the resulting run log is eligible for
%     ingestion once the agent completes the task.
% \end{description}
%
% Signals compose: the executor applies signals from all active detectors in
% sequence before dispatching each action.  A single deviation may trigger
% \texttt{refresh\_observation} followed by \texttt{run\_actions} in the same
% step.  Each detector is an independent plugin; adding or removing a detector
% does not change the executor or the Function format.  The set of active
% detectors is configurable per deployment.
